\documentclass[12pt]{exam}
\pagestyle{headandfoot}
\firstpageheader{}{}{}
\runningheader{\date}{}{\class, \teacher}
\runningheadrule
\firstpagefooter{}{}{}
\runningfooter{}{Viel Erfolg!}{Seite \thepage\:von \numpages}
\runningfootrule
%=========================================================
\newcommand{\theme}{Rationale Zahlen}
\newcommand{\teacher}{Philipp Moor}
\newcommand{\class}{2b}
\newcommand{\serie}{A}
\newcommand{\duration}{45 Minuten}
\newcommand{\name}{\bf Name:}
\newcommand{\grade}{Note:}
\newcommand{\datum}{22.09.22}
%=========================================================

%=========================================================
\pagestyle{headandfoot}
\firstpageheader{}{}{}
\runningheader{\datum}{}{Name: \hspace{5cm}}
\runningheadrule
\firstpagefooter{}{}{}
\runningfooter{}{Viel Erfolg!}{Seite \thepage\ von \numpages}
\runningfootrule
%========================================================
\begin{document}
\fontsize{14}{14}\selectfont
%========================================================

%========================================================
\begin{tabular*}{\textwidth}{l @{\extracolsep{\fill}}l @{\extracolsep{6pt}}l}
\textbf{\theme} & \textbf{\grade} & {\answerbox}\\
\textbf{Klasse \class} & & \\
\textbf{\teacher}  & {\bf\name} &\textbf{\hrulefill}\\[8pt]
%-------------------
\end{tabular*}
%--------------------
\begin{center}
\rule[1ex]{\textwidth}{1pt}
Dauer: \duration \hspace{9cm}Datum: \datum\\
\vspace{5pt}
\rule[2ex]{\textwidth}{1pt}
\end{center}
%=========================================================

%=========================================================
\begin{center}
\textbf{Bewertungstabelle}\\
\small Wird von der Lehrperson ausgefüllt\\
\vspace{5pt}
\addpoints
\gradetable[h][questions]\\
\vspace{5pt}
Taschenrechner und jegliche andere Hilfsmittel sind \textbf{nicht erlaubt}.\\
Schreibe alle Blätter mit \textbf{Vor}-und \textbf{Nachnamen} an.
\end{center}

\noindent
\rule[1ex]{\textwidth}{1pt}

\vspace{10cm}
\begin{center}
    \Huge{Serie \serie}
\end{center}

\newpage

\begin{questions}
%=========================================================
%-------------------TESTFRAGEN
%=========================================================
\question[3]
Rechne aus, welchen Anteil des folgenden $7\times3$ - Rechtecks die rot eingefärbte Fläche ausmacht. Schreibe das Resultat als Bruch und kürze soweit wie möglich.

\newpage
%---------------------------------------------------------
\question[6]
Kürze die folgenden Brüche soweit wie möglich. Markiere bei deinen Schritten mittels durchstreichen welche Terme du miteinander gekürzt hast.
\begin{parts}
\setlength{\columnsep}{30pt}
\begin{multicols}{3}
\part $\frac{195}{-45}$
\part $\frac{-24z^4}{-120x^2z^{8}}$
%\part $\frac{12x^2y^4}{72x^2y^{12}}$
\part $\frac{121x^6y}{77x^6}$
\end{multicols}
\end{parts}
\vspace{15pt}
\antwort{20}

\newpage
\question[6]
Berechne jeweils für (a) und (b) die Lücken ($\triangle,\square,\bigcirc$):
\begin{parts}
\setlength{\columnsep}{30pt}
\part $\frac{3}{7} = \frac{\triangle}{28} = \frac{3^2}{\square} = \frac{\bigcirc}{-42xy}$
% = \frac{45z^4}{\star}$
\vspace{5pt}
\antwort{11}
\part $6x = \frac{\triangle}{-2} = \frac{18x^2}{\square} = \frac{\bigcirc}{30x^3}$
%= \frac{72xy^5}{\star}$
\vspace{5pt}
\antwort{9}
\end{parts}
\vspace{15pt}


\newpage
\question[6]
Berechne die folgenden Summen und Differenzen und kürze das Resultat soweit wie möglich:
\begin{parts}
\setlength{\columnsep}{30pt}
\part $\frac{1}{2}-\left[ \frac{3}{4}-\left( \frac{5}{8}+\frac{2}{10}\right)\right]$
\part $\frac{15}{315}-\frac{x}{42x}-\frac{5}{30}+\frac{2}{14x}$
\end{parts}
\vspace{15pt}
\antwort{20}

\newpage
\question[6]
Berechne das Resultat und kürze es soweit wie möglich:
\begin{parts}
\setlength{\columnsep}{30pt}
\part $\left(\frac{24xy}{33}:\frac{12x^2}{132}\right)-\frac{16y}{2x}$
\part $ \left[3 - \left(\frac{12}{2x}\cdot\frac{4}{6}\right):\frac{4}{x}\right]:\frac{1}{4} $
\end{parts}
\vspace{15pt}
\antwort{20}

\newpage
\question[6]
Berechne x und kürze das Resultat so weit wie möglich:
\begin{parts}
\setlength{\columnsep}{30pt}
\part $\frac{8}{20}\cdot x = 2$
\part $x \cdot \frac{4}{3} = 10$
\part $x \cdot \left(\frac{1}{2}\right)^2 = 1$
\end{parts}
\vspace{15pt}
\antwort{20}

\addpoints
\end{questions}
\end{document}